\chapter{Why Do Revolutions Happen?}
\section{Paper Worth Less and Less}
Pick up a piece of paper. More precisely, a banknote.

A little larger than your palm, its rough paper carries elaborate ornament, solemn emblems and lines of French. At its centre stands the denomination: fifty livres. A Paris mason working long days might earn twenty or thirty livres a month. When newly printed, this sheet therefore represented nearly two months of an able-bodied worker's sweat.

It has a special name: assignat. France's revolutionary government began issuing these in late 1789. Their story still sounds modern. Revolution left the treasury empty and government owing astronomical debts. The revolutionaries boldly nationalised Catholic Church land and issued notes secured against it. The reasoning sounded solid: real land stood behind every note, redeemable through land purchases. They should be as reliable as gold.

Initially, people believed. Then war came, spending burst its banks, and notes were printed far faster than land sold. More assignats meant less value. By 1796, the revolution's seventh year, they traded for a fraction of face value; baskets bought little bread. Memoirs describe daily panic: a housewife set out with enough for four pounds of bread, hesitated an afternoon and could afford only three. Eventually government invalidated them. A revolution opening with liberty, equality and fraternity ended in ordinary wallets as comprehensive inflation.

Study the paper longer. It layers all our questions: how could national finances collapse until printed paper sustained them? Why did people proclaiming equality and liberty first confiscate Church property? Why did revolution overshoot its original destination into war, terror and dictatorship? Most fundamentally, how does a society-upending revolution happen? Was it predetermined, or produced by colliding accidents?

\section{Paris in July 1789}
Come into a scene.

The time is late on 13 July into the early hours of 14 July 1789. Paris holds over six hundred thousand people, among Europe's largest cities, with magnificent palaces and sewage-filled alleys alike.

Remember bread. Working households spent over half their income on it: bread was life. France's 1788 harvest suffered summer drought and extraordinary hail just before reaping, followed by a bitter winter freezing the Seine. By summer 1789, bread prices reached levels rarely seen for decades. Before dawn on 14 July, queues formed outside neighbourhood bakeries. People gripped worn copper coins, calculating whether bread would rise again.

Another rumour frightened the city more: royal troops were massing outside Paris. To dissolve the new National Assembly? To massacre the city? Nobody knew. Those words were perfect panic fuel. Days earlier, the king had suddenly dismissed the trusted finance minister Necker. Many read this as a signal that he was about to act.

On the morning of 14 July, crowds stormed the Invalides and seized tens of thousands of muskets. Guns needed powder. Where was it? The Bastille.

The eight-towered medieval fortress stood in eastern Paris. Popular imagination made it royal power's dark symbol, holding prisoners arrested by secret orders. In fact, only seven remained that day, none quite the imagined political prisoners. Symbols do not audit inventories. Commander de Launay's hundred-odd defenders faced thousands of agitated citizens. Negotiations failed, shots sounded and fighting continued for hours in narrow streets. Towards evening, the garrison surrendered. Nearly a hundred attackers died, mainly local locksmiths, carpenters and small shopkeepers. De Launay was dragged out and killed, his head paraded on a pike. The seven prisoners were released and carried in procession.

News left Paris that night, crossed France in days and Europe in weeks. A falling stone fortress sounded like a falling age.

Keep two details: the attackers were not the most destitute, but skilled workers and small proprietors. Half their impetus was bread, half rumour, many July rumours later proving false.

\section{The Question Is Not Who Was Right}
The textbook question waits: what was the French Revolution's historical significance? Ignore it briefly. We pursue a clumsier, harder question:

\textbf{Why do revolutions happen?}

It invites glib false answers. A wicked king? Far worse monarchs than Louis XVI died peacefully. Suffering people? History contains worse suffering in places without revolution. Radical ideas? Genevans and Scots reading the same Voltaire and Rousseau did not guillotine their monarchs. Move each smooth answer elsewhere and it fails.

That same summer across the Atlantic, roughly half a million enslaved people laboured without freedom in France's richest colony, Saint-Domingue, now Haiti, supplying Europe with sugar and coffee. Months later, Paris's National Assembly would solemnly proclaim everyone born free and equal while carefully avoiding slavery. Colonial wealth supported French finances and many deputies' fortunes.

Our question therefore has two layers. Explanatorily, which combination of finance, society, ideas and timing exploded the old order? Ethically, how should we judge the gap between proclaimed ideals and actual conduct? Both are necessary and must remain separate. Explaining a revolution does not approve it; condemning violence does not make its words meaningless.

\section{Different Accounts inside and Outside}
Revolution is not good people against bad. On the eve of collapse, different positions revealed different ledgers. Hear each fully, beginning with the least sympathetic: the king and his finance ministers. This is steelmanning, strengthening an opponent's case before testing it rather than searching for weaknesses.

\textbf{The crown: we are repaying debts, not squandering money.}

From Versailles's accounts, enormous public debt included a major bill for helping American independence. France paid money and sent fleets to free Britain's colonies, then received the invoice at home. Interest consumed roughly half government revenue. Repayment required taxes, but nearly every route was blocked. Nobles and clergy held extensive land with exemptions, burdening the Third Estate, especially peasants. Challenge privilege and the aristocratic judges of the Parlement of Paris refused to register new tax laws. This parlement was not a modern trial court, but a political institution reviewing royal commands. Reforms from Turgot through Necker to Calonne crashed against that wall.

Little here is false. Louis XVI was no cartoon tyrant: he commissioned reforms, agreed to summon Estates-General dormant for 175 years and later swore to uphold a constitution. From his position, removing privileges required more than a signature. The state was built on estates; nobody could guarantee removing one brick would not collapse its walls. Present this fully: the old order died not because nobody wanted repairs, but because willing repairers could not make them.

\textbf{The Third Estate: we support the nation yet count for nothing.}

In January 1789, the Abbé Sieyès's pamphlet What Is the Third Estate? swept Paris. Its three answers were, broadly: what is it? Everything. What has it been politically? Nothing. What does it seek? To become something. Everyone outside clergy and nobility belonged to it, over ninety per cent of France: bankers, lawyers, peasants and porters. They farmed, traded, paid taxes and fought, without adequate national influence. The pamphlet ignited France not through intricate theory, but by printing countless people's resentment: why did household supporters lack a voice while non-taxpayers held a veto?

\textbf{Saint-Domingue: does your everyone include us?}

Across the ocean, the eighteenth century's most profitable colony enriched French merchants through cane and coffee at the cost of some half a million enslaved lives. Working lives were frighteningly short; planters calculated that working people to death and replacing them cost less than improving conditions. Alongside were about thirty thousand whites and twenty to thirty thousand free people of colour. Revolutionary news gave each group its own reading of liberty and equality. Small white planters wanted political rights, free people of colour equality with whites, and slaves asked whether everyone born free included them. In August 1791, northern slaves rose and plantations burned in rows at night. A former enslaved coachman, Toussaint Louverture, later assumed leadership.

\textbf{The planters: your declaration calls property sacred and inviolable.}

This was no word game: Article 17 of the Declaration of the Rights of Man and of the Citizen explicitly protected property, while colonial law made slaves property. On the planters' reading, emancipation violated the revolution's own declaration. This logical knot would entangle it; we shall revisit the text.

\textbf{An almost uninvited voice: women.}

In 1791, Paris playwright Olympe de Gouges modelled a Declaration of the Rights of Woman and of the Female Citizen article by article on the earlier document. Men were born free and equal: what about women? Two years later, she was guillotined. Revolution gave her the declaration's syntax without intending to give her a place within it.

Together the five voices show something crucial. Every party could find supporting revolutionary principles. The tearing came not from reasonable people opposing unreasonable ones, but from liberty, equality, property and security colliding without yielding.

\section[Thinking Bridge: Four Layers of Why]{Thinking Bridge: Divide ``Why'' into Four Layers}
The worst answer to revolution's cause gives one factor; the next worst lists many. A useful tool separates why into four layers answering distinct questions. Imagine a forest fire:

\textbf{First: deep structure---what is the forest made of?} Slow variables over decades or centuries include taxation constrained by privilege, rigid estates and population-driven food-price pressures. They determine fire-prone places, not ignition dates. A wooden house standing safely for a century does not make wood harmless.

\textbf{Second: medium-term circumstances---how dry have recent seasons been?} Over several years come American-war debts, Enlightenment books and pamphlets, and intensifying friction between nobles and rising property owners. Drought raises danger without igniting itself.

\textbf{Third: triggers---the spark.} Events over days or weeks include the 1788 natural disasters, Estates-General, Necker's dismissal and troop rumours. Sparks fall everywhere; on wet ground they do nothing, in dry resinous pine they bring disaster.

\textbf{Fourth: choices---some extinguish, others fan.} At critical points, alternatives existed. The king could concede to the Third Estate in June instead of hardening his stance. Parisians could wait rather than attack the Bastille. Nobles could defend privilege on 4 August instead of dramatically renouncing it. Revolution is no stone rolling downhill: decision-makers stand before every steep descent.

The tool prevents two shortcuts. Structural determinism calls revolution inevitable, absolving every 1789 decision. Yet favourable weather in 1788 might have let the Estates-General become peaceful constitutional reform, like England's 1688 elite-negotiated transfer with little bloodshed at home. Accidental determinism makes everything luck, unable to explain why bankrupt, estate-divided France exploded rather than neighbouring Geneva.

Try China in 1911. Deep structure: indemnities crushing finances and scholars displaced by abolition of examinations. Medium-term conditions: constitutionalists' complete disillusionment with court reform. Triggers: Sichuan's Railway Protection Movement and an accidental discharge in Wuchang. Choices: provincial assembly leaders joining revolution and Yuan Shikai declining further service to the dynasty. These four layers reveal a much smaller fire than France's: the wealthiest and strongest chose containment over fanning, at the cost of leaving most old accounts unsettled.

\section[Two Declarations and Their Cracks]{Reading Evidence: Two Declarations and Their Cracks}
Revolutionary years left mountains of writing. Two declarations, American and French, reward close reading.

On 4 July 1776 in Philadelphia, thirteen colonies' representatives adopted the Declaration of Independence, opening as follows in the common Chinese translation rendered here:

\begin{quote}
We regard these truths as self-evident: all people are created equal, endowed by their Creator with certain inalienable rights, including life, liberty and the pursuit of happiness.
\end{quote}
---The American Declaration of Independence, 1776.

Its principal drafter was Thomas Jefferson. Remember him; we shall return shortly.

On 26 August 1789, France's National Assembly issued the Declaration of the Rights of Man and of the Citizen. Article 1 reads:

\begin{quote}
People are born and remain free and equal in rights.
\end{quote}
Article 17 of the same document reads:

\begin{quote}
Property is a sacred, inviolable right. Nobody may be deprived of it unless legally established public necessity clearly requires this, with fair compensation paid in advance.
\end{quote}
---The French Declaration of the Rights of Man and of the Citizen, 1789, Articles 1 and 17.

Set them side by side. Freedom and equality in Article 1; sacred property in Article 17. In metropolitan France they could coexist. In Saint-Domingue they collided because colonial law defined slaves as property. Strictly implementing Article 1 meant freeing half a million immediately; strictly implementing Article 17 meant leaving owners' property untouched. The declaration did not specify precedence. People had to answer through negotiation, votes and finally fire and blades.

Answers came rapidly and harshly. The August 1791 uprising burned over a thousand northern-plain plantations. French officials proclaimed local emancipation in 1793; the National Convention abolished slavery across all French colonies in 1794, the first European great power to legislate abolition. Glory did not end the story. In 1802, Napoleon sent an expedition to restore the old order and reinstated slavery elsewhere. Toussaint was deceived into going to France and died in 1803 in a prison near the Alps. His comrades won the war and declared Haiti independent on 1 January 1804, the modern world's first state founded by a slave uprising.

Return to Jefferson. The author of equality owned hundreds of enslaved people over his lifetime at his Virginia estate and freed very few at death. This does not invalidate the declaration: calling Jefferson hypocritical and calling equality false are different claims. Many historians prefer the image of a cheque drawn on the future. Its writer did not intend payment, but could not control it after circulation. Haitian slaves, American abolitionists and countless later movements presented it for redemption. Revolutionary texts strangely draw strength from precisely what their authors did not believe.

The fissure spread beyond Haiti. At dawn on 16 September 1810, Father Hidalgo rang the church bell in Dolores, Mexico, igniting Indigenous and mixed-race peasants' anger. The Cry of Dolores remains a national-day memory. Defeated and executed the next year, he passed on the flame. In South America, Bolívar's and San Martín's armies defeated Spanish colonial rule over more than a decade. By the 1820s, Spain's American holdings were virtually reduced to Caribbean islands, with new republics across the mainland. Outcomes differed, however. Most movements were led by creoles, American-born Spanish-descended elites seeking transfer from European masters, not social inversion. Some republics abolished slavery quickly; others retained it for decades. Brazil's independence from Portugal in 1822 kept even an emperor, with slavery surviving until 1888. Under the same independence slogan, some revolutions changed society, others only flags. Remember that difference.

\section{Three Accounts of the French Revolution}
Why France in 1789? Two centuries of scholarly argument make comparing explanations more useful than simply accepting one. We compare causes, not the evidential fact that the Bastille fell on 14 July.

\textbf{Explanation one: fiscal crisis---bankruptcy opened the door.}

This account studies government books. Royal debts accumulated, privilege riddled taxation, and borrowing cost more than Britain's. By the late 1780s, even interest payments faltered. Summoning the Estates-General resembled a bankrupt company convening creditors. This solid account explains timing and an overlooked fact: the state fired revolution's opening shot itself, gathering national elites and giving them a bargaining platform. Sociologically, old power created a political opportunity.

Its limitation is equally solid: fiscal crises are common, revolutions rare. Britain offers a counterexample, with greater per-capita public debt but no revolution. Through Parliament, taxpayers participated in decisions; taxes bought influence, parliamentary backing made debt credible and borrowing cheaper. France collapsed not from debt alone but from the discredited arrangement of payment without voice. Fiscal explanations alone resemble matches explaining fires without explaining wooden houses.

\textbf{Explanation two: social structure---the gears seized.}

Nobles defended exemptions and court incomes; rising lawyers, merchants and rich people had money without access; peasants endured layered feudal dues. Three rusted gears ground against one another. Structural accounts explain the energy: such a fire required generations of accumulated resentment.

Beware two traps. First, the cartoon bourgeois revolution has long been revised. Deputies were mainly lawyers and officials; radicals included nobles such as Lafayette and Mirabeau; Parisian crowds and castle-burning peasants pursued their own interests, not a single class marching uniformly. Second, structure explains disorder, not its particular shape. Similar structures yielded parliamentary negotiation in England in 1688, Terror and Napoleon in France in 1789. People held the steering wheel at forks.

\textbf{Explanation three: circulating ideas---minds changed before the world.}

Voltaire and Rousseau circulated, salons discussed reason, pamphlets attacked despotism, and American independence demonstrated that popular sovereignty could build a state. This explains why 1789 was not merely another tax riot, rapidly generating a new language of nation, rights and constitution. Without it, storming the Bastille meant stealing powder; with it, popular insurrection. Tocqueville added a sharper observation in The Old Regime and the Revolution: revolutions often begin not under the worst oppression, but amid improvement. Before 1789, reform made remaining injustice less tolerable. A bad government's most dangerous moment may be when it starts improving.

Yet most peasants burning deeds had not read Rousseau. Attributing half a million people's actions to a few dozen books is a storyteller's romance. Ideas resemble an ammunition store, not a trigger: they do not explain ignition's timing, but the name under which people burn.

The accounts divide labour: finance explains timing, structure energy, ideas justification. Forcing them into one correct answer is wrong; treating them as equivalent is also wrong. For the specific year 1789, finance explains more than ideas. Comparison means asking what each explains and cannot explain, not indiscriminate compromise.

\section[Further Challenge: Situations and Outcomes]{Further Challenge: Revolutionary Situations Are Not Revolutionary Outcomes}
Charles Tilly, an American sociologist who spent his career studying revolution and collective action, distinguishes \textbf{revolutionary situations} from \textbf{revolutionary outcomes}.

A revolutionary situation is rare: competing centres claim the same population's allegiance and old authority cannot contain them. Tilly calls it multiple sovereignty. Recall France in 1789: Versailles's crown, Paris's National Assembly and the Paris Commune in the streets simultaneously issuing orders. In China in 1911, Wuchang's military government and Beijing's Qing court each claimed the nation. Russia in 1917 shows it more plainly: Provisional Government and Petrograd Soviet coexisted facing one another for over half a year. The core is not anger, which occurs daily, but a cracked ceiling of power allowing allegiance to switch teams.

Outcomes are different: who fills the crack? Organisation, weapons, money and legitimacy decide. Tilly's cold observation is that revolutionary situations are commoner than revolutionary outcomes. Often negotiation, repression or reform closes the crack, or its occupant differs entirely from insurgents' wishes. Through this lens our cases reorder:

England in 1688 transferred power through elite negotiation and an invitation, leaving little time for dual power and little domestic bloodshed. Do not romanticise it: the same revolution brought years of Irish war. France's dual power pulled for three years, yielding Jacobin dictatorship and the Terror of 1793--1794. Revolutionary tribunals formally executed about seventeen thousand; with civil wars such as the Vendée, tens of thousands died. Napoleon's imperial crown followed. Haiti's deepest crack destroyed old power through war; insurgents' own armies filled it. Its revolution became the most thorough, changing rulers and removing slavery itself from law. Across much of Latin America, creole elites controlled the process, preserving hierarchy beneath new flags.

Tilly illuminates an uncomfortable truth: \textbf{ideals and anger do not determine revolutionary endings; organisational capacity does}. The best slogans seldom guarantee the final victory. This is not cynicism, but a measuring rod. When news shows rival power centres, look beyond slogans to the distribution of guns, money, organisation and legitimacy.

\section[Bread, Rumours and Revolution Today]{Today: Bread Prices, Rumours and the Word ``Revolution''}
These two-hundred-year-old mechanisms are nearby.

In December 2010, Tunisian street vendor Mohamed Bouazizi's cart was confiscated. Unable to secure redress, he set himself alight outside a government building. Within weeks, protests spread nationally and President Ben Ali fled after twenty-three years in power. Review the event and old acquaintances appear: rising food prices, youth unemployment's structural resentment, phone videos spreading one person's death nationwide within hours, and members of the ruling coalition refusing further shooting for old power. Bread, dry fuel, upgraded rumour and human choices: four interlocking layers, no sole fuse.

A quieter echo appears in constitutions worldwide. Preambles about born freedom and popular sovereignty inherit formulations from 1776 and 1789. Revolution's longest-lived legacy is repeated sentences, not barricades. Its other legacy remains too. Twenty-one years after Haitian independence, in 1825, French warships demanded compensation for former planters' property losses under threat of war. Haiti assumed an independence debt, repaid intermittently for over a century. The world's unique successful slave-revolt state was fined for its masters' losses. Outcomes depend on post-war debt, blockade and recognition as well as battles.

Notice also revolution's light modern usage: industrial, information and AI revolutions, revolutionary product upgrades. Usually it is marketing without anyone seeking to overthrow power. Language can move; that is not necessarily wrong. But distinguish the weight: one revolution changes tools, another stakes lives. When someone casually proposes revolution as a cure, ask four questions: what is the fuel, what is the spark, have you counted the Terror's cost, and will the person filling the crack be the one you want?

\section{Your Turn: Who Should Settle the Account?}
Return to the assignat.

Solemn revolutionary promises printed on one face became waste paper. It miniaturises revolution: between ideals and fulfilment stand presses, guillotines, wars and long debts.

Consider a final question without a standard answer. Imagine being born enslaved in Saint-Domingue in 1791. Two paths lie ahead. Wait for Parisian consciences and gradual reform: later history showed a wait of over a decade, interrupted by rulers like Napoleon, with mother and children spending their lives in cane fields. Or burn: join the uprising, set plantations alight and take up a blade. History also showed massacre, revenge, victorious generals becoming dictators and humiliating compensation in 1825 along this path.

\textbf{Obedience costs your whole life; crossing the boundary costs innocent blood. Which do you choose?}

Do not answer too quickly. If waiting, fully present the insurgents' case before rejecting it. If burning, fully present ordinary people caught in war before maintaining your choice. Then set out evidence: the declarations, collapsed finance, fractured power and its eventual replacement.

Push further. Many terms we use to judge past revolutions---liberty, equality and rights---were invented or developed by those revolutions. Is measuring revolution with its own ruler circular? Could another measure, such as a cane-field child's lifespan, produce a different result?

There is no standard answer. But you have been calculating since picking up the note. Revolution is not confined to museums. It is an arithmetic problem returning to humanity in different clothes, requiring every generation to decide how to settle its account.
